\chapter{Bosonic control with ancilla and coupler}
\label{preliminary}
\section{Oscillator Circuit Hamiltonian}
\label{sec:oscillator_hamiltonian}

The simplest quantized electrical mode is the lossless LC oscillator, which serves as the foundational circuit element of superconducting quantum circuits. Its importance is twofold. First, it provides a direct correspondence between circuit variables and canonical variables in quantum mechanics. Second, it illustrates a fundamental limitation: a purely linear oscillator, even when driven by classical fields, cannot generate non-Gaussian states such as Fock states from the vacuum. This limitation motivates the introduction of a nonlinear element in the next subsection.

\subsection{Quantization of the LC oscillator}

We begin with an inductor $L$ and a capacitor $C$ connected in parallel. The natural generalized coordinate is the node flux
\begin{equation}
    \Phi(t) = \int_{-\infty}^{t} V(t')\,dt',
\end{equation}
where $V(t)$ is the node voltage. The conjugate variable is the node charge $Q$, which is associated with the capacitor branch. In terms of the flux coordinate, the classical Lagrangian is
\begin{equation}
    \mathcal{L}(\Phi,\dot{\Phi}) = \frac{C}{2}\dot{\Phi}^{2} - \frac{1}{2L}\Phi^{2}.
\end{equation}
The first term is the electrostatic energy stored in the capacitor, while the second term is the magnetic energy stored in the inductor.

The canonical momentum conjugate to $\Phi$ is
\begin{equation}
    Q \equiv \frac{\partial \mathcal{L}}{\partial \dot{\Phi}} = C\dot{\Phi}.
\end{equation}
This is precisely the capacitor charge. Performing the Legendre transform gives the Hamiltonian
\begin{equation}
    H = Q\dot{\Phi} - \mathcal{L}
      = \frac{Q^{2}}{2C} + \frac{\Phi^{2}}{2L}.
\end{equation}
Thus, the LC circuit is formally identical to a mechanical harmonic oscillator, with the correspondence
\begin{equation}
    \Phi \leftrightarrow x,
    \qquad
    Q \leftrightarrow p,
    \qquad
    C \leftrightarrow m,
    \qquad
    \frac{1}{L} \leftrightarrow m\omega^{2}.
\end{equation}
The resonance frequency is
\begin{equation}
    \omega_{0} = \frac{1}{\sqrt{LC}}.
\end{equation}

Canonical quantization promotes $\Phi$ and $Q$ to operators satisfying
\begin{equation}
    [\hat{\Phi},\hat{Q}] = i\hbar.
\end{equation}
The quantum Hamiltonian is therefore
\begin{equation}
    \hat{H}_{0} = \frac{\hat{Q}^{2}}{2C} + \frac{\hat{\Phi}^{2}}{2L}.
\end{equation}
It is convenient to introduce the zero-point fluctuation amplitudes
\begin{equation}
    \Phi_{\mathrm{zpf}} = \sqrt{\frac{\hbar Z}{2}},
    \qquad
    Q_{\mathrm{zpf}} = \sqrt{\frac{\hbar}{2Z}},
\end{equation}
where
\begin{equation}
    Z = \sqrt{\frac{L}{C}}
\end{equation}
is the characteristic impedance of the oscillator. We then define annihilation and creation operators through
\begin{equation}
    \hat{\Phi} = \Phi_{\mathrm{zpf}}(\hat{a}+\hat{a}^{\dagger}),
    \qquad
    \hat{Q} = -iQ_{\mathrm{zpf}}(\hat{a}-\hat{a}^{\dagger}),
\end{equation}
with $[\hat{a},\hat{a}^{\dagger}] = 1$. Substituting these expressions into $\hat{H}_{0}$ yields
\begin{equation}
    \hat{H}_{0}
    =
    \hbar\omega_{0}\left(\hat{a}^{\dagger}\hat{a}+\frac{1}{2}\right),
\end{equation}
which is the standard quantum harmonic oscillator Hamiltonian.

The eigenstates of this Hamiltonian are the Fock states $\{\ket{n}\}_{n=0}^{\infty}$, defined by
\begin{equation}
    \hat{a}^{\dagger}\hat{a}\ket{n} = n\ket{n}.
\end{equation}
These states have equally spaced energies,
\begin{equation}
    E_{n} = \hbar\omega_{0}\left(n+\frac{1}{2}\right).
\end{equation}
Because the level spacing is uniform, the linear oscillator does not distinguish one transition from another. This spectral degeneracy is the central obstacle to selective state preparation using only linear control.

\subsection{Classical driving of the oscillator}

The oscillator can be controlled by applying a classical microwave tone. For a capacitive drive, the control enters linearly through the charge or flux quadrature. A convenient form for the driven Hamiltonian is
\begin{equation}
    \hat{H}(t)
    =
    \hat{H}_{0}
    -
    \epsilon(t)\hat{\Phi},
\end{equation}
where $\epsilon(t)$ is a classical drive amplitude determined by the external voltage source and coupling geometry. Expressed in ladder operators, this becomes
\begin{equation}
    \hat{H}(t)
    =
    \hbar\omega_{0}\hat{a}^{\dagger}\hat{a}
    +
    f(t)\hat{a}
    +
    f^{*}(t)\hat{a}^{\dagger},
\end{equation}
up to an irrelevant constant, where $f(t)$ is proportional to the classical microwave envelope.

For a monochromatic tone near resonance,
\begin{equation}
    f(t) = \hbar\Omega e^{-i\omega_{d} t},
\end{equation}
the Hamiltonian reads
\begin{equation}
    \hat{H}(t)
    =
    \hbar\omega_{0}\hat{a}^{\dagger}\hat{a}
    +
    \hbar\Omega e^{-i\omega_{d} t}\hat{a}
    +
    \hbar\Omega^{*} e^{i\omega_{d} t}\hat{a}^{\dagger}.
\end{equation}
Moving to a frame rotating at the drive frequency and applying the rotating-wave approximation gives
\begin{equation}
    \hat{H}_{\mathrm{rot}}
    =
    \hbar\Delta\,\hat{a}^{\dagger}\hat{a}
    +
    \hbar\Omega\,\hat{a}
    +
    \hbar\Omega^{*}\hat{a}^{\dagger},
\end{equation}
where $\Delta = \omega_{0}-\omega_{d}$ is the detuning.

The corresponding time-evolution operator is a displacement operator,
\begin{equation}
    \hat{U}(t) = e^{-i\phi(t)} \hat{D}(\alpha(t)),
\end{equation}
where
\begin{equation}
    \hat{D}(\alpha) = \exp\!\left(\alpha\hat{a}^{\dagger}-\alpha^{*}\hat{a}\right).
\end{equation}
Therefore, a linearly driven harmonic oscillator maps the vacuum state to a coherent state,
\begin{equation}
    \ket{\psi(t)} = \hat{D}(\alpha)\ket{0} = \ket{\alpha}.
\end{equation}
More generally, if the Hamiltonian remains quadratic in $\hat{a}$ and $\hat{a}^{\dagger}$, classical controls can generate only Gaussian unitaries. Besides displacement, this broader class includes squeezing through terms of the form
\begin{equation}
    \hat{H}_{\mathrm{sq}}(t)
    =
    \hbar\lambda(t)\hat{a}^{2}
    +
    \hbar\lambda^{*}(t)\hat{a}^{\dagger 2},
\end{equation}
which arise from parametric modulation rather than from a simple linear microwave drive. The associated unitary is the squeezing operator
\begin{equation}
    \hat{S}(\xi)
    =
    \exp\!\left[
        \frac{1}{2}
        \left(
            \xi^{*}\hat{a}^{2}
            -
            \xi \hat{a}^{\dagger 2}
        \right)
    \right].
\end{equation}
Hence, the full reachable set of a classically controlled \emph{linear} oscillator is restricted to Gaussian transformations: displacements, phase-space rotations, and, when quadratic time-dependent control is allowed, squeezing.

\subsection{Why Fock states are unreachable in a linear oscillator}

This restriction has an immediate consequence for state preparation. The vacuum state $\ket{0}$ is Gaussian, and Gaussian unitaries preserve Gaussianity. Therefore, any state obtained from $\ket{0}$ under a Hamiltonian that is at most quadratic in $\hat{a}$ and $\hat{a}^{\dagger}$ must also be Gaussian. Symbolically,
\begin{equation}
    \ket{\psi_{\mathrm{reachable}}}
    =
    \hat{U}_{\mathrm{Gaussian}}\ket{0}
    \quad \Longrightarrow \quad
    \ket{\psi_{\mathrm{reachable}}}
    \text{ is Gaussian.}
\end{equation}

Fock states with $n\geq 1$ are non-Gaussian. For example, their Wigner functions exhibit oscillations and negativity, which cannot be produced by Gaussian evolution starting from vacuum. In particular, no choice of classical linear drive can generate $\ket{1}$, $\ket{2}$, or any other excited number state from $\ket{0}$. The obstruction is not merely practical but structural: the Hamiltonian of a harmonic system preserves the Gaussian character of the state and lacks the nonlinearity required to isolate specific transitions in the equally spaced ladder.

Equivalently, one may view the problem from the spectrum. Because every adjacent transition has the same energy $\hbar\omega_{0}$, a resonant drive that couples $\ket{0}\leftrightarrow\ket{1}$ also couples $\ket{1}\leftrightarrow\ket{2}$, $\ket{2}\leftrightarrow\ket{3}$, and so on. The drive cannot spectrally select a single pair of levels. As a result, linear control produces phase-space motion rather than deterministic preparation of a specific non-Gaussian Fock state.

This limitation provides the essential motivation for introducing a nonlinear circuit element. By breaking the exact harmonicity of the spectrum, a nonlinear oscillator acquires nonuniform level spacing, which enables selective addressing of individual transitions and allows the generation of genuinely non-Gaussian quantum states. In superconducting circuits, the most important realization of such a weakly anharmonic oscillator is the transmon, which we introduce in the next subsection.